\documentclass[10pt,a4paper]{amsart}
\usepackage[margin=19mm]{geometry}
\usepackage{amsmath,amssymb,mathtools}
\usepackage[scr=rsfs,cal=euler]{mathalfa}
\usepackage{xfrac}
\usepackage[unicode,hidelinks,bookmarks=false]{hyperref}
\newcommand{\ie}{i.e.\ }
\newcommand{\cf}{cf.\ }
\newcommand{\resp}{respectively\ }
\newcommand{\Subsection}{\subsection}
\providecommand{\nobreakdash}{\nobreak\hbox{-}}
\setlength{\emergencystretch}{2em}
\multlinegap0pt
\usepackage{booktabs}
\usepackage{mathtools}
\usepackage{xcolor}
\usepackage{amssymb}
\usepackage{caption}
\usepackage{siunitx}
\usepackage{float}
\usepackage{dsfont}
\usepackage{enumerate}
\usepackage[shortlabels]{enumitem}
\usepackage{braket}
\usepackage{tikz}
\newtheorem{theorem}{Theorem}
\newcommand{\Id}{\mathrm{Id}}
\newcommand{\R}{\mathbb{R}}
\newcommand{\di}{\, \mathrm{d}}
\newcommand{\rhobar}{\overline{\rho}}
\newcommand{\sph}{\mathbb{S}}
\newcommand{\tra}{\mathrm{tr}}
\title{A new infinitesimal form of the Pr\'ekopa--Leindler inequality with multiplicative structure and applications}
\author{Sotiris Armeniakos, Jacopo Ulivelli}
\date{Source: arXiv:2602.10812v2 (CC BY 4.0)}
\begin{document}
\maketitle

\noindent\emph{Source and licence.} A new infinitesimal form of the Pr\'ekopa--Leindler inequality with multiplicative structure and applications. Version arXiv:2602.10812v2 (CC BY 4.0), distributed by arXiv under the Creative Commons Attribution 4.0 International licence, http://creativecommons.org/licenses/by/4.0/, as recorded in the arXiv licence metadata for this exact version and shown on the abstract page https://arxiv.org/abs/2602.10812v2. Full licence notice: \texttt{licenses/arxiv-2602.10812-CC-BY-4.0.txt}.\par\smallskip

\noindent\textit{Editorial scope.} Counted unit: the article's theorem thm:reformulation (source0.tex lines 907--913). The article's complete original proof of it is delivered with it (source0.tex lines 914--981, one \texttt{proof} environment of 68 lines). No mathematics is written, repaired or re-derived: the statement and the proof are the article's own lines, in the article's own notation, and no retained line is substituted or re-flowed.  Retained uncounted as the source-local context the proof needs: lines 194--200; lines 836--838; lines 890--892.  Nothing was omitted from any retained span, so the package carries no omission record.  No retained line contains a citation, so the wrapper carries no bibliography and no bibliography entry.  No retained line contains a figure environment, an image-inclusion command or drawing code, so no image is delivered and the package declares no asset dependency.  Editorial formatting only, in addition: an AMS \texttt{amsart} preamble that restates the article's own declarations and macros used by the retained lines, namely the environments \texttt{theorem} on separate counters, together with the standard packages those lines need.  The retained environments print in this document as Theorem 1 in order of appearance, whereas the article prints the same items with the numbering of its own full document.  The complete native source of the article as distributed in the arXiv e-print for this exact version is bundled unchanged as \texttt{sources/194430/source0.tex}.
\begin{theorem}\label{thm:multiplicative_form}
    Consider a non-empty compact convex set $K \subset \R^n$ such that $\partial K$ is a strictly convex manifold of class $C^2$, and a convex function $u \in C^2(\R^n)$ such that its Hessian matrix is positive definite. Then for every Lipschitz function $\rho: \partial K \to \R$ and every locally Lipschitz function $\varphi: \R^n \to \R$ the inequality 
    \begin{equation}\label{eq:multiplicative_form}
        \langle \rho, \varphi \rangle_{\rm I}^2 \leq \langle \rho, \rho \rangle_{\rm P} \langle \varphi, \varphi \rangle_{\rm BL}
    \end{equation}
    holds.
\end{theorem}

\begin{equation}\label{eq:new_concavity expression}
p(\mu,K) = \frac{\mu(K)}{\int_{\partial K} \overline{\rho}\,\di\mu}.
\end{equation}

\begin{equation}\label{eq:weighted_div}
    \int_{\partial K} \nabla_{\partial K, \mu} \cdot X \di \mu = 0
\end{equation}

\begin{theorem}\label{thm:reformulation}
    Let $K$ and $\mu$ satisfy the assumptions of Theorem~\ref{thm:multiplicative_form}. In particular, $\di \mu(x)=e^{-u(x)}\di x$ for a smooth and strictly convex function $u$. Then, 
    \begin{equation}\label{eq:equivalent_dim_BM}
        p(\mu,K) \geq \frac{1}{n} \Longleftrightarrow \langle \rhobar,\langle \nabla u, x \rangle \rangle_{\rm I}+\int_K \langle \nabla u, x \rangle \di \mu \geq 0,
    \end{equation}
    where $\rhobar$ is the unique solution of $\mathcal{L}(\rho)=1$ on $\partial K$.
\end{theorem}

\begin{proof}
We start by observing that 
\begin{equation}\label{eq:integral_support}
    \int_{\partial K} h_{K}(\nu_{K})\,\di \mu
= n\,\mu(K) -\int_{K}\langle \nabla u,x\rangle \,\di \mu.
\end{equation}
Similar statements have appeared in the literature. We provide a proof for the convenience of the reader. Indeed, for $x\in \partial K$ one has $h_{K}(\nu_{K}(x))=\langle x,\nu_{K}(x)\rangle$. The divergence theorem yields
\begin{align*}
&\int_{\partial K}h_{K}(\nu_{K})\,\di\mu
=\int_{\partial K} \langle x, \nu_{K}\rangle e^{-u}\,\di x=
\int_{K} \nabla \cdot \left(x e^{-u}\right)\,\di x \\
=&\int_{K} n\,\di\mu -\int_{K}\langle \nabla u,x\rangle\,\di \mu 
=n\mu(K)-\int_{K}\langle \nabla u(x),x\rangle\,\di \mu,
\end{align*}
proving \eqref{eq:integral_support}.

Next, we claim that
\begin{equation}\label{eq:support_calculation}
\mathcal{L}\big(h_{K}(\nu_{K})\big) = 1+\langle \nabla u,x\rangle -\frac{1}{\mu(K)}\int_{K} \langle \nabla u, x\rangle \,\di \mu.
\end{equation}
First, observe that, by the properties of the support function
\begin{align*}
    \nabla_{\partial K} \cdot \left( {\rm II}^{-1} \nabla_{\partial K}(h_K(\nu_K))\right)= & \nabla_{\partial K}\cdot \left( \nabla_{\sph^{n-1}}h_K(\nu_K)\right)\\
    = \tra\left( \Id -h_K(\nu_K){\rm II}\right)= &(n-1)-h_K(\nu_K)\tra({\rm II}).
\end{align*}
Therefore,
\begin{equation}\label{eq:weighted_part}
\nabla_{\partial K, \mu} \cdot \left({\rm II}^{-1}\nabla_{\partial K}\big(h_{K}(\nu_{K})\big)\right) =(n-1)-h_K(\nu_K)\tra({\rm II})
-\langle \nabla_{\partial K}u, x\rangle .
\end{equation}
Combining \eqref{eq:weighted_part} with the definition of $\mathcal{L}$, the definition of ${\rm H}_\mu$, and \eqref{eq:integral_support}, we obtain
\begin{align*}
\mathcal{L}\big(h_{K}(\nu_{K}(x))\big)
&=
1+\langle \nabla_{\partial K}u(x), x\rangle
+\langle \nabla u(x), \nu_{K}(x)\rangle\,h_{K}(\nu_{K}(x))
-\frac{1}{\mu(K)} \int_{K} \langle \nabla u(x), x\rangle \,\di \mu(x) \\
&=
1+\langle \nabla u(x), x\rangle
-\frac{1}{\mu(K)}\int_{K} \langle \nabla u(x), x\rangle \,\di \mu(x),
\end{align*}
and \eqref{eq:support_calculation} holds. 

Observe now that, analogously to \eqref{eq:weighted_div}, the operator $\mathcal{L}$ is symmetric in the sense that for any sufficiently regular functions $\rho,\sigma: \partial K \to \R$,
\[
\int_{\partial K} \rho\mathcal{L}(\sigma)\di \mu
=
\int_{\partial K} \mathcal{L}(\rho)\,\sigma\,\di \mu.
\]
Using this property together with the fact that $\mathcal{L}(\overline{\rho})=1$ in the weak sense, we obtain
\begin{align*}
\int_{\partial K}h_{K}(\nu_{K})\,\di\mu
-
\int_{\partial K}\overline{\rho}\,\di \mu
=
\int_{\partial K} \big(h_{K}(\nu_{K})-\overline{\rho}\big)\,\mathcal{L}(\overline{\rho})\,\di \mu 
=
\int_{\partial K} \mathcal{L}\big(h_{K}(\nu_{K})-\overline{\rho}\big)\,\overline{\rho}\,\di \mu.
\end{align*}
A direct application of \eqref{eq:support_calculation} yields
\begin{equation} \label{eq:thm_dim_ref}
    \int_{\partial K}h_{K}(\nu_{K})\,\di\mu- \int_{\partial K} \overline{\rho}\,\di \mu = \langle \overline{\rho}, \langle \nabla u , x \rangle \rangle_{\rm I}.
\end{equation}

Finally, by \eqref{eq:new_concavity expression} and \eqref{eq:integral_support}, $p(\mu,K)\geq 1/n$ is equivalent to
\[n\mu(K) \geq \int_{\partial K}\rhobar \di \mu  \Leftrightarrow \int_{\partial K} h_K(\nu_K) \di \mu +\int_{K} \langle \nabla u,x \rangle \di \mu \geq \int_{\partial K} \rhobar \di \mu,\]
which, in turn, by \eqref{eq:thm_dim_ref} is equivalent to \[\langle \rhobar,\langle \nabla u, x \rangle \rangle_{\rm I}+\int_K \langle \nabla u, x \rangle \di \mu \geq 0,\] concluding the proof.
\end{proof}

\end{document}
